% !TEX root = ../PhD.tex

\section{SDS}\label{sec:sds}

\noindent\textit{Multiple datasets can be derived from the same data, by
transforming the objects on the incoming data stream. We want to give query
processors transparency in how these datasets are related and what they
contain. In this section, we introduce the Smart Data Specification for
Semantically Describing Streams (SDS) to annotate datasets with provenance
information, describing the consumed stream and the applied transformations
on that stream. We demo the execution of a pipeline that transforms a Linked
Data Event Stream and publishes the data in a different structure as
described in the SDS description. The SDS vocabulary and application profile
bring together DCAT-AP, LDES and P-Plan.}
\vspace{1em}

\subsection{Introduction}

Data portals publish live datasets derived from streams. Streams are often
related to other streams by transforming their data. This leads to problems
when datasets are not accompanied by provenance information: the data portal
has difficulty knowing whether or not this data is already published. Also,
during dataset selection, query agents are clueless as to what interfaces
provide the same underlying data and, as a fallback, they query the same
data multiple times.

For example, in a GPS application that notifies the user about changes in
routes due to construction sites, multiple (query) interfaces can be used. A
time-based interface makes it easy to track the latest changes. On the other
hand, a geospatial-based interface makes it easy to calculate whether or not
a construction site will be encountered on a route. A SPARQL endpoint can
fulfil both interfaces, but SPARQL carries high operation costs and may
leave the users with low availability~\cite{SPARQL}. Linked Data Event
Streams (LDES) alleviate these high costs by publishing static web resources
that can be easily cached~\cite{LDES}. Each stream can be fragmented in a
different way to accommodate the needs of the users.

Tailoring dataset interfaces to particular needs results in a plethora of
interfaces, which makes well-defined provenance a necessity. The provenance
should cover the two steps a query agent takes to query an
interface~\cite{QueryDiscovery1,QueryDiscovery2}:

\begin{enumerate}
  \item \emph{Dataset discovery}: based on whether a dataset is going to
    contain statements that contribute to solving the query
  \item \emph{Interface discovery}: selecting the right interface that
    publishes the dataset
\end{enumerate}

We set up a new interface for the Belgian street name registry and
demonstrate that a query agent can find the optimal interface to execute
particular queries.

\subsection{Related Work}

\textbf{DCTerms, DCAT and VoID:} Exposing metadata about datasets is long
established. Dublin Core Terms (DCTerms) can be used to provide basic
information about resources, providing terms like \emph{title},
\emph{author}, \emph{description} and \emph{subject}~\cite{DCTerms}. Data
Catalog Vocabulary (DCAT) is designed to facilitate interoperability between
data catalogs published on the web~\cite{DCAT}. DCAT also provides terms
like \emph{license}, which makes it possible to define a new license for an
interface. The Vocabulary of Interlinked Datasets (VoID) focuses on
explicitly linking datasets together on some predicate and defining subset
datasets~\cite{VOID}.

\textbf{LDES:} Linked Data Event Streams is a way of exposing an
ever-growing set of immutable objects. These objects can be divided into
fragments as HTTP resources that are linked together with the TREE
specification. Fragmentations are used to provide semantic meaning to links
between HTTP resources. Each HTTP resource can, for example, hold all items
that start with a particular letter. A view description is used to define
the meaning of the fragments and their links~\cite{LDES}.

\textbf{VoCaLS:} Vocabulary of interoperable streams \& On a Web of Data
Streams extends the ideas of DCAT with more information about streaming
data~\cite{VoCaLS}. That work defines a stream slightly differently than in
this section. VoCaLS focuses on streams that generate high throughput
updates, which requires processors to use a windowing mechanism. Here, a
stream is seen more broadly as a growing collection of objects, updates or
otherwise.

\textbf{P-Plan and PROV-O:} The Ontology for Provenance and Plans (P-Plan)
is an extension of the PROV-O ontology~\cite{PROVO} created to represent the
plans that guide the execution of scientific processes. P-Plan describes how
plans are composed. This information can be used to keep track of a
provenance tree~\cite{PPLAN}.

\subsection{The Smart Data Specification for Semantically Describing
Streams (SDS)}

A stream in the context of the Smart Data Specification for Semantically
Describing Streams (SDS) is a \emph{physical} live channel that carries
updates or items. A dataset can be derived from a stream as the collection
of all updates or items. A \emph{physical} channel can be any medium such as
a Kafka stream, a WebSocket stream, or even a file where updates are
appended. A stream can carry any kind of data: CSV rows, mutable or
immutable linked data objects, video stream chunks, etc.

A stream can be derived from a transformation applied to items of a
different stream. This transformation is described with \texttt{p-plan} and
resides in the SDS description of the resulting stream. The stream and the
transformation correspond with \texttt{p-plan:Entity} and
\texttt{p-plan:Activity} respectively, as shown in the pink part of
Figure~\ref{fig:sds-onto}. The transformation links to the previous stream
with the \texttt{prov:used} predicate. With the power of \texttt{p-plan},
query agents can understand how datasets are linked and what interface fits
a specific query the best.

The SDS description also links to metadata about the resulting data with
the \texttt{sds:dataset} predicate. This makes modifying the datasets'
metadata possible after a transformation, represented as the green part in
Figure~\ref{fig:sds-onto}.

A \texttt{sds:Stream} should describe the type of objects it is carrying.
First, the stream can use the \texttt{sds:carries} predicate that points to
the type of the \texttt{sds:Records} it is carrying (see below). Next, the
stream can describe the SHACL shape with \texttt{sds:shape} that these
members adhere to, which can only be used if the stream carries linked data
objects.

Linking specific items to the correct stream is done with
\texttt{sds:Record} objects. An \texttt{sds:Record} points to the data
(\texttt{sds:payload}) and the corresponding stream (\texttt{sds:stream}).
These small objects make it possible for multiple streams to use the same
channel. Each transformation can thus push \texttt{sds:Record} objects and
leave the original stream intact. A stream of immutable objects can still be
transformed, for example, to calculate a hash or add a fragment id to the
\texttt{sds:Record} object. The yellow part of Figure~\ref{fig:sds-onto}
gives a visual overview of \texttt{sds:Record}.

\begin{figure}[!ht]
  \centering
  \includegraphics[width=0.7\textwidth]{sds_ontology.png}
  \caption{Visual representation of the SDS Ontology. The colored boxes
    have no semantic meaning and are only used for clarity.}
  \label{fig:sds-onto}
\end{figure}

\subsection{Demo}

Data published with Linked Data Event Streams can be partitioned or
fragmented in a multitude of ways. This helps query agents resolve their
queries as fast as possible whilst ingesting as little data as possible. A
default fragmentation constitutes a timestamp fragmentation, which allows
clients to replicate and synchronize the dataset efficiently. A substring
fragmentation, on the other hand, makes autocompletion more
efficient~\cite{Substring}.

In this demo, we set up a pipeline starting from an existing LDES that
exposes the registry of street names with a timestamp fragmentation. The
pipeline calculates a substring fragmentation based on the name of the
street and exposes a new LDES with the corresponding SDS description and
substring fragmentation. The published \texttt{tree:View}s can then be
associated with their respective \texttt{viewDescription}, as shown in
Listing~\ref{lst:sds-example}.

\begin{lstlistingsp}[language=ttl,caption={RDF sample creating two LDES Views from different Streams},label={lst:sds-example},basicstyle=\ttfamily\scriptsize]
@prefix ex:     <http://example.org/ns#>.
@prefix rdfs:   <http://www.w3.org/2000/01/rdf-schema#> .
@prefix ldes:   <https://w3id.org/ldes#>.
@prefix p-plan: <http://purl.org/net/p-plan#> .
@prefix prov:   <http://www.w3.org/ns/prov#> .
@prefix sds:    <https://w3id.org/sds#> .

ex:MyLDES a ldes:EventStream;
  dcat:title "An example LDES".

# One fragmentation is based on an existing LDES
ex:BasicFragmentation a tree:ViewDescription ;
  dcat:endpointURL </basic> ; # A rootnote from which you can access all members
  dcat:servesDataset ex:MyLDES ; # the LDES
  # This viewDescription was created by importing this stream
  ldes:managedBy sds:LDESStream.

# The other fragmentation bucketizes the members before publishing
ex:SubstringFragmentation a tree:ViewDescription ;
  dcat:endpointURL </substring> ; # A rootnote from which you can access all members
  dcat:servesDataset ex:MyLDES ; # the LDES
  # This viewDescription was created by importing this stream
  ldes:managedBy sds:BucketizedStream.


ex:ImportLDES a p-plan:Activity;
  rdfs:comment "Reads csv file and converts to rdf members";
  prov:used <https://smartdata.dev-vlaanderen.be/base/gemeente>.

ex:LDESStream a sds:Stream;
  p-plan:wasGeneratedBy ex:ImportLDES;
  sds:carries sds:Member.


ex:BucketizeStrategy a ldes:BucketizeStrategy;
    ldes:bucketType ldes:SubstringFragmentation;
    tree:path rdfs:label;
    ldes:pageSize 50.

ex:BucketizeStream a p-plan:Activiy;
  rdfs:comment "Execute a substring bucketization on the incoming stream";
  prov:used ex:LDESStream, ex:BucketizeStrategy.

ex:BucketizedStream a sds:Stream;
  p-plan:wasGeneratedBy ex:BucketizeStream;
  sds:carries sds:Member.
\end{lstlistingsp}

When asking a query agent ``What are the 10 latest updated street names?''
starting from the newly created LDES, the query agent can derive from the
SDS description that the current LDES is not suitable for this query. This
query would require the query agent to request the entire LDES tree and
manually find the 10 latest updates, whereas following the links from the
SDS description back to the original LDES, this query would only require a
few HTTP requests: one to get the SDS description and another to get the
latest updates due to the timestamp-based fragmentation.

To execute this pipeline we use a proof-of-concept pipeline runner called
Nautirust~\cite{Nautirust}. This makes it easy to start the three required
processes with the correct arguments: read the original LDES with an LDES
client, add buckets to the SDS Records, and ingest the new SDS records in an
LDES server.

\subsection{Conclusion}

With the introduction of the SDS ontology, it is possible to add a
description to a stream and the resulting dataset, adding provenance
information. The provenance links streams together and the transformations
applied to them. The SDS ontology aligns well with long-established
ontologies like DCAT and P-Plan to maximize interoperability.

With the SDS description, the query agent can now automatically select the
right dataset and interface based on a given query.

Federated query processors that utilize source selection based on this
provenance information when selecting a dataset and interface to query were
still future work at the time of the original 2022 demo paper.

\subsection{Revisiting SDS Against the Current State of Play}

\todo{Compare the SDS vocabulary (2022) against how LDES/DCAT-AP/dataset-provenance
tooling has evolved since; note what has changed in the ecosystem, whether
the federated source-selection future work materialized, and how SDS relates
to MuMo's own use of LDES and provenance today.}
